\section{Appendix: Energy-Based AC Resistance Analysis}\label{sec:appendix_ac_resistance}

\subsection{Discrete AC Equation near Curved Walls}

The Allen-Cahn equation in the phase-field LBM framework is discretized as:
\begin{equation}\label{eq:app_ac}
    \phi^{n+1} = \phi^n + \Delta t \left[ M \nabla^2 \phi - \frac{M}{\xi^2} \phi(1-\phi)(1-2\phi) + \nabla \cdot (\phi \mathbf{u}) \right]
\end{equation}
Near a curved solid wall, the geometric wetting BC modifies the gradient $\nabla\phi$ at the first fluid node. The AC sharpening source term in the normal direction is:
\begin{equation}\label{eq:app_ac_source}
    S_{AC,n} = \frac{M}{\xi^2} \phi(1-\phi)(1-2\phi) \cdot \mathbf{n}
\end{equation}

At the interface center ($\phi = 0.5$), the nonlinear term $\phi(1-\phi)(1-2\phi)$ vanishes, but the gradient $\nabla\phi$ is maximum. The net AC resistance arises from the \emph{gradient of the source term} rather than the source term itself.

\subsection{Asymptotic Analysis}

Consider a one-dimensional interface profile near a flat wall at $x=0$. The steady-state AC equation gives:
\begin{equation}\label{eq:app_ac_1d}
    M \frac{d^2\phi}{dx^2} = \frac{M}{\xi^2} \phi(1-\phi)(1-2\phi)
\end{equation}
The exact solution is $\phi(x) = \frac{1}{2}(1 + \tanh(x/\xi))$, with gradient $d\phi/dx = 1/(2\xi \cosh^2(x/\xi))$. At the interface center ($x=0$), the normal gradient is $\left.d\phi/dx\right|_{x=0} = 1/(2\xi)$. The AC source term at $x=0$ is zero (since $\phi=0.5$), but the \emph{curvature} of $\phi$ creates a restoring force that acts as a spring resisting deviations from the equilibrium profile.

\subsection{Resistance Force Derivation}

The AC sharpening resistance can be quantified by considering the energy barrier for deforming the interface. The free energy density is:
\begin{equation}\label{eq:app_free_energy}
    f(\phi) = \frac{\beta}{2} \phi^2(1-\phi)^2 + \frac{\kappa}{2} |\nabla\phi|^2
\end{equation}

For a geometric wetting correction that shifts the interface position by $\delta$ near the wall, the energy cost is $\Delta f \approx \kappa |\nabla\phi|^2 \delta^2 / \xi$. The resistance force (energy gradient) is $R_{AC} = \partial f/\partial \delta \approx \kappa |\nabla\phi|^2 \delta / \xi$. Using $\kappa = \beta\xi^2/8$ and $|\nabla\phi| \approx 1/(2\xi)$ at the interface: $R_{AC} \approx \beta\delta/(32\xi)$. The target gradient from the geometric wetting condition scales as $|g_n^{target}| \sim 2/\xi$. For strongly hydrophobic surfaces ($\theta_{eq} \gg 90^\circ$), the required correction $\delta$ is large, making $R_{AC}$ significant relative to $g_n^{target}$.

\subsection{Physical Interpretation}

The free energy functional for the Allen-Cahn model~\cite{allen1979microscopic, cahn1958free} is:
\begin{equation}\label{eq:app_free_energy_full}
    F[\phi] = \int \left[ \frac{\beta_{FE}}{2} \phi^2(1-\phi)^2 + \frac{\kappa}{2} |\nabla\phi|^2 \right] dV
\end{equation}
where $\beta_{FE} = 12\sigma/\xi$ is the bulk free energy parameter and $\kappa = \beta_{FE}\xi^2/8$ is the interface stiffness~\cite{jacqmin2000contact, lowengrub1998quasi}. For a geometric wetting correction that shifts the interface position by $\delta$ near the wall, the energy cost comes primarily from the gradient term. For a tanh interface profile, the gradient energy density is $\frac{\kappa}{8\xi^2} \text{sech}^4(x/\xi)$.

Integrating over the interface region yields the energy change $\Delta F \approx \frac{\kappa}{6\xi} \delta^2 G(\delta/\xi)$, where $G$ is a geometric correction factor accounting for nonlinear profile deformation. This energy argument provides a qualitative mechanism for the AC resistance; a quantitative prediction would require the dynamic interface state during impact, which is not modeled here.

The scaling law for the amplification factor follows from this analysis:
The amplification magnitude follows the main-text scaling law (Eq.~\ref{eq:alpha_min}), $\alpha_{geo} = 1 + C(\xi/R_{\mathrm{eff}})\,|\cot\theta_{eq}|$ with $C = 0.914$ calibrated at $R^*=1.0$, $\theta_{eq}=162^\circ$ (see the Supplementary Material).

\subsection{Sensitivity Analysis}

 Qualitatively, the strongest sensitivity is to interface thickness $\xi$, which directly affects the AC resistance length scale.
